% Options for packages loaded elsewhere
\PassOptionsToPackage{unicode}{hyperref}
\PassOptionsToPackage{hyphens}{url}
\PassOptionsToPackage{space}{xeCJK}
\documentclass[
  a4paper,
  ja=minimal,
  xelatex,
  paragraph-mark={},
]{bxjsarticle}
\usepackage{xcolor}
\geometry{margin=25mm}
\usepackage{amsmath,amssymb}
\usepackage{mathrsfs}
\setcounter{secnumdepth}{-\maxdimen} % remove section numbering
\usepackage{iftex}
\ifPDFTeX
  \usepackage[T1]{fontenc}
  \usepackage[utf8]{inputenc}
  \usepackage{textcomp} % provide euro and other symbols
\else % if luatex or xetex
  \usepackage{unicode-math} % this also loads fontspec
  \defaultfontfeatures{Scale=MatchLowercase}
  \defaultfontfeatures[\rmfamily]{Ligatures=TeX,Scale=1}
\fi
\usepackage{lmodern}
\ifPDFTeX\else
  % xetex/luatex font selection
  \ifXeTeX
    \usepackage{xeCJK}
    \setCJKmainfont[]{Hiragino Mincho ProN}
    \setCJKsansfont[]{Hiragino Kaku Gothic ProN}
    \setCJKmonofont[]{Hiragino Kaku Gothic ProN}
  \fi
  \ifLuaTeX
    \usepackage[]{luatexja-fontspec}
    \setmainjfont[]{Hiragino Mincho ProN}
  \fi
\fi
% Use upquote if available, for straight quotes in verbatim environments
\IfFileExists{upquote.sty}{\usepackage{upquote}}{}
\IfFileExists{microtype.sty}{% use microtype if available
  \usepackage[]{microtype}
  \UseMicrotypeSet[protrusion]{basicmath} % disable protrusion for tt fonts
}{}
\makeatletter
\@ifundefined{KOMAClassName}{% if non-KOMA class
  \IfFileExists{parskip.sty}{%
    \usepackage{parskip}
  }{% else
    \setlength{\parindent}{0pt}
    \setlength{\parskip}{6pt plus 2pt minus 1pt}}
}{% if KOMA class
  \KOMAoptions{parskip=half}}
\makeatother
\usepackage{longtable,booktabs,array}
\newcounter{none} % for unnumbered tables
\usepackage{calc} % for calculating minipage widths
% Correct order of tables after \paragraph or \subparagraph
\usepackage{etoolbox}
\makeatletter
\patchcmd\longtable{\par}{\if@noskipsec\mbox{}\fi\par}{}{}
\makeatother
% Allow footnotes in longtable head/foot
\IfFileExists{footnotehyper.sty}{\usepackage{footnotehyper}}{\usepackage{footnote}}
\makesavenoteenv{longtable}
\setlength{\emergencystretch}{3em} % prevent overfull lines
\providecommand{\tightlist}{%
  \setlength{\itemsep}{0pt}\setlength{\parskip}{0pt}}
\usepackage{bookmark}
\IfFileExists{xurl.sty}{\usepackage{xurl}}{} % add URL line breaks if available
\urlstyle{same}
\hypersetup{
  hidelinks,
  pdfcreator={LaTeX via pandoc}}

\author{}
\date{}

\begin{document}

\section{無名部分投射による物理的対称性選択の公理的枠組み}\label{ux7121ux540dux90e8ux5206ux6295ux5c04ux306bux3088ux308bux7269ux7406ux7684ux5bfeux79f0ux6027ux9078ux629eux306eux516cux7406ux7684ux67a0ux7d44ux307f}

\subsection{\texorpdfstring{八成分二クァルテット分枝における \(E_8\)
の条件付き同定}{八成分二クァルテット分枝における E\_8 の条件付き同定}}\label{ux516bux6210ux5206ux4e8cux30afux30a1ux30ebux30c6ux30c3ux30c8ux5206ux679dux306bux304aux3051ux308b-e_8-ux306eux6761ux4ef6ux4ed8ux304dux540cux5b9a}

\textbf{著者:} 木原 範昭\\
\textbf{ORCID:} 0009-0004-6753-4020\\
\textbf{日付:} 2026年7月24日\\
\textbf{Version DOI:}
\href{https://doi.org/10.5281/zenodo.21521900}{10.5281/zenodo.21521900}\\
\textbf{Concept DOI:}
\href{https://doi.org/10.5281/zenodo.21521899}{10.5281/zenodo.21521899}\\
\textbf{位置づけ:} 独立論文・日本語版 v1／数理物理学・物理学基礎論

\begin{center}\rule{0.5\linewidth}{0.5pt}\end{center}

\subsection{要旨}\label{ux8981ux65e8}

基礎公理の解集合が複数の閉鎖系を含むとき、観測される物理的対称性を解集合全体の対称性と同一視する必要はない。本稿は、物理的状態空間を、基礎公理が許す候補空間から選ばれた部分投射像として定義する\textbf{無名部分投射選択原理}を提示する。

候補状態空間を
\(\mathscr S\)、無名性・スケール無名性・閉鎖整合性・読出し整合性を満たす許容部分投射の集合を
\(\mathfrak P(\mathscr S)\) とする。本稿の中心公理は、

\[
\boxed{
\exists\mathcal P_*\in\mathfrak P(\mathscr S)
\quad\text{such that}\quad
\mathscr S_{\mathrm{phys}}
{}=
\operatorname{Im}\mathcal P_*
}
\]

である。この公理は、どの対称群を選ぶかを指定しない。物理的対称性は、選択された像に備わる関係構造の自己同型群として、投射後に同定される。

本枠組みの非自明な適用例として、選択像の遷移差が、正定値八次元読出し空間内で二つの四成分系に分かれ、各局所遷移格子が
\(D_4\) 型となる分枝を扱う。二つの \(D_4\)
判別群に同一の中心作用を課し、その中性成分だけを残す共通中心射影を採ると、許容接着類は

\[
(0,0),\qquad(v,v),\qquad(s,s),\qquad(c,c)
\]

に限られる。この対角接着で得られる格子は、正定値・八次元・偶・ユニモジュラーであり、したがって
\(E_8\) 格子と同型である。最短ベクトルは

\[
48+3\times64=240
\]

本となり、対応する複素単純 Lie 代数の次元は

\[
8+240=248
\]

となる。さらに \(E_8\) の Coxeter 元は位数30を持つ。

本稿は、基礎公理が \(E_8\)
を唯一解として強制すると主張しない。八成分性、二クァルテット \(D_4\)
構造、共通中心射影が物理的に選ばれる理由は、本稿では選択問題として残る。本稿が確立するのは、無名な部分投射選択を先に定式化し、その特定分枝を数学的に分類した結果として
\(E_8\) が現れる、という公理的順序と条件付き構造定理である。

\begin{center}\rule{0.5\linewidth}{0.5pt}\end{center}

\subsection{0.
位置づけとスコープ}\label{ux4f4dux7f6eux3065ux3051ux3068ux30b9ux30b3ux30fcux30d7}

\subsubsection{0.1
本稿が扱うもの}\label{ux672cux7a3fux304cux6271ux3046ux3082ux306e}

本稿は次の三点を扱う。

\begin{enumerate}
\def\labelenumi{\arabic{enumi}.}
\tightlist
\item
  基礎公理の許容解空間と物理的実現空間を分離する。
\item
  物理的実現を無名な許容部分投射の像として定式化する。
\item
  八成分二クァルテット分枝に共通中心中性射影を課した場合、選択像の遷移格子が
  \(E_8\) になることを証明する。
\end{enumerate}

\subsubsection{0.2
本稿が扱わないもの}\label{ux672cux7a3fux304cux6271ux308fux306aux3044ux3082ux306e}

本稿は次を主張しない。

\begin{enumerate}
\def\labelenumi{\arabic{enumi}.}
\tightlist
\item
  基礎公理だけから \(E_8\) が一意に導出されること。
\item
  すべての物理的部分投射像が \(E_8\) 対称性を持つこと。
\item
  実際の投射 \(\mathcal P_*\)
  を選ぶ力学または選択汎関数が既に得られていること。
\item
  八つの読出し量と、粒子・相互作用・標準模型の各量との対応が確立していること。
\item
  有限回帰作用素が物理的に \(E_8\) の Coxeter 元と同一であること。
\item
  \(E_8\) が時空全体または全状態空間の対称性であること。
\end{enumerate}

\subsubsection{0.3 中心主張}\label{ux4e2dux5fc3ux4e3bux5f35}

本稿の中心主張は、次の順序である。

\[
\boxed{
\begin{aligned}
&\text{基礎公理}\\
&\Longrightarrow
\text{複数の許容閉鎖系からなる候補空間}\\
&\Longrightarrow
\text{無名な許容部分投射の集合}\\
&\Longrightarrow
\text{一つの物理的部分像の選択}\\
&\Longrightarrow
\text{選択像の対称性の数学的同定}.
\end{aligned}
}
\]

\(E_8\)
はこの鎖の最初に置かれる公理ではなく、特定分枝における最後の同定結果である。

\begin{center}\rule{0.5\linewidth}{0.5pt}\end{center}

\subsection{1.
基礎公理との接続}\label{ux57faux790eux516cux7406ux3068ux306eux63a5ux7d9a}

\subsubsection{1.1
継承する形成規則}\label{ux7d99ux627fux3059ux308bux5f62ux6210ux898fux5247}

本稿は、無名等振幅複合波モデル基本公理系
v8［1］から、次の形成規則を継承する。

\textbf{無名性:}
基本成分に、個別名、特権軸、特権符号、特権型を与えない。物理名または既存理論名を理由として、第一原理層の式または読出し規則を変更しない。

\textbf{スケール無名性:} 状態 \(X\) と
\(\lambda X\)、\(\lambda\in\mathbb C^\ast\)
は、絶対スケールを除いた関係構造として同値である。

\textbf{全正符号ゼロ閉鎖:}

\[
Q(X)
:=
\sum_{j=1}^{N}x_j^2
{}=0.
\]

\textbf{非自明存在:}

\[
X\neq0.
\]

\subsubsection{1.2
本稿で追加するもの}\label{ux672cux7a3fux3067ux8ffdux52a0ux3059ux308bux3082ux306e}

上の形成規則は候補状態を制限するが、候補のうち何が物理的に読まれるかを一意に指定しない。本稿は、この空白を埋めるために、\textbf{部分投射の存在と許容性}を定式化する。

本稿は特定の対称群を追加公理として採用しない。特に、

\[
\text{「物理的対称性は }E_8\text{ である」}
\]

を公理として置かない。これは、求める理論名を理由に読出し規則を選ぶことを禁じる無名性と両立しないからである。

\begin{center}\rule{0.5\linewidth}{0.5pt}\end{center}

\subsection{2. 許容解空間}\label{ux8a31ux5bb9ux89e3ux7a7aux9593}

\subsubsection{2.1
射影化された候補状態}\label{ux5c04ux5f71ux5316ux3055ux308cux305fux5019ux88dcux72b6ux614b}

スケール同値関係を、

\[
X\sim_{\mathrm{sc}}\lambda X,
\qquad
\lambda\in\mathbb C^\ast
\]

とする。非自明零閉鎖候補の射影化された集合を、

\[
\mathscr S_0
:=
\left\{
[X]
\ \middle|\
X\neq0,\ Q(X)=0
\right\},
\]

\[
[X]
:=
\left\{
\lambda X
\mid
\lambda\in\mathbb C^\ast
\right\}
\]

と定義する。

\subsubsection{2.2
発展を含む候補系}\label{ux767aux5c55ux3092ux542bux3080ux5019ux88dcux7cfb}

状態空間上の複素線形発展作用素を \(U\) とする。線形性により、\(U\)
はスケール同値類上にも作用する。同じ閉鎖系の発展であるための適合条件を、

\[
Q(UY)=Q(Y)
\qquad
(\forall Y)
\]

とする。

発展を含む候補系の集合を、

\[
\mathscr S
:=
\left\{
([X],U)
\ \middle|\
[X]\in\mathscr S_0,\ Q\circ U=Q
\right\}
\]

と定義する。

\subsubsection{2.3
有限回帰分枝}\label{ux6709ux9650ux56deux5e30ux5206ux679d}

有限回帰を持つ候補系を、

\[
\mathscr S_{\mathrm{fin}}
:=
\left\{
([X],U,n)
\ \middle|\
([X],U)\in\mathscr S,
\qquad
U^n=I,
\qquad
n\in\mathbb N
\right\}
\]

とする。

有限回帰は本稿の一般的な部分投射原理には必須ではない。これは、後に
\(E_8\) の Coxeter 周期と接続可能な部分族を指定する追加条件である。

\begin{center}\rule{0.5\linewidth}{0.5pt}\end{center}

\subsection{3. 無名部分投射}\label{ux7121ux540dux90e8ux5206ux6295ux5c04}

\subsubsection{3.1 定義}\label{ux5b9aux7fa9}

\textbf{定義 3.1（部分投射）:} 候補空間 \(\mathscr S\) の部分集合
\(\mathscr D_{\mathcal P}\subseteq\mathscr S\) と写像

\[
p_{\mathcal P}:
\mathscr D_{\mathcal P}
\longrightarrow
\mathscr Y_{\mathcal P}
\]

の組

\[
\mathcal P
{}=
\left(
\mathscr D_{\mathcal P},
p_{\mathcal P}
\right)
\]

を部分投射と呼ぶ。

ここで部分投射は、線形代数における直交射影に限定しない。次の三操作を含み得る。

\begin{enumerate}
\def\labelenumi{\arabic{enumi}.}
\tightlist
\item
  候補系の一部だけを定義域として残す選択。
\item
  読出し不能な差異を同一視する商。
\item
  読出し可能な関係量を別の表示空間へ写す写像。
\end{enumerate}

\textbf{定義 3.2（非自明部分投射）:}
次の少なくとも一方を満たす部分投射を非自明と呼ぶ。

\[
\mathscr D_{\mathcal P}
\subsetneq
\mathscr S,
\]

または、

\[
\exists s_1\neq s_2
\quad\text{such that}\quad
p_{\mathcal P}(s_1)
{}=
p_{\mathcal P}(s_2).
\]

すなわち、候補の一部が除かれるか、候補間の差異の一部が読出し不能として同一視される。恒等写像だけを置く場合は、本稿の部分投射選択原理の非自明な実現とは数えない。

\subsubsection{3.2 遷移差格子}\label{ux9077ux79fbux5deeux683cux5b50}

選択像に離散的な遷移差が存在するとする。選択された遷移差集合を
\(\Delta_{\mathcal P}\subseteq V_{\mathcal P}\) とし、その整数包を、

\[
\Lambda_{\mathcal P}
:=
\operatorname{span}_{\mathbb Z}
\Delta_{\mathcal P}
\]

と定義する。

\(\Lambda_{\mathcal P}\)
は、状態そのものの個別名を格納するのではなく、読出し可能な状態差または遷移差を格納する。

\subsubsection{3.3 許容性条件}\label{ux8a31ux5bb9ux6027ux6761ux4ef6}

\textbf{定義 3.3（許容部分投射）:}
非自明であり、かつ次の条件を満たす部分投射を許容部分投射と呼ぶ。

\paragraph{A. 成分無名性}\label{a.-ux6210ux5206ux7121ux540dux6027}

成分の置換または無名性を保存する変換 \(g\)
に対し、投射結果は個別ラベルの付け替えに依存しない。

\[
p_{\mathcal P}(g\cdot[X])
\sim_{\mathscr Y}
p_{\mathcal P}([X]).
\]

\paragraph{B.
スケール無名性}\label{b.-ux30b9ux30b1ux30fcux30ebux7121ux540dux6027}

\[
p_{\mathcal P}([\lambda X])
{}=
p_{\mathcal P}([X]),
\qquad
\lambda\in\mathbb C^\ast.
\]

\paragraph{C. 閉鎖整合性}\label{c.-ux9589ux9396ux6574ux5408ux6027}

投射の定義域は基礎閉鎖を満たす候補から構成され、投射後の遷移差は、投射像に定義された合成則のもとで閉じる。

\paragraph{D.
読出し整合性}\label{d.-ux8aadux51faux3057ux6574ux5408ux6027}

投射によって同一視された二状態は、採用された読出し量だけからは区別されない。

\paragraph{E.
物理名先行の禁止}\label{e.-ux7269ux7406ux540dux5148ux884cux306eux7981ux6b62}

投射の定義に、投射後にのみ与えられる物理名または対称群名を用いない。

許容部分投射全体を、

\[
\mathfrak P(\mathscr S)
\]

と書く。

\begin{center}\rule{0.5\linewidth}{0.5pt}\end{center}

\subsection{4.
部分投射存在公理}\label{ux90e8ux5206ux6295ux5c04ux5b58ux5728ux516cux7406}

\subsubsection{4.1 公理}\label{ux516cux7406}

\textbf{公理 PP（無名部分投射存在公理）:}

基礎公理の候補空間 \(\mathscr S\)
に対し、物理的に実現される状態空間は、少なくとも一つの非自明な許容部分投射の像である。

\[
\boxed{
\exists\mathcal P_*
\in
\mathfrak P(\mathscr S)
\quad\text{such that}\quad
\mathscr S_{\mathrm{phys}}
{}=
\operatorname{Im}\mathcal P_*.
}
\]

\subsubsection{4.2
公理が主張しないこと}\label{ux516cux7406ux304cux4e3bux5f35ux3057ux306aux3044ux3053ux3068}

公理 PP は、

\[
\mathcal P_*
\text{ が一意である}
\]

とは主張しない。また、

\[
\mathcal P_*
{}=
\operatorname*{arg\,max}_{\mathcal P\in\mathfrak P(\mathscr S)}
\mathcal F(\mathcal P)
\]

のような選択汎関数 \(\mathcal F\) の具体形も指定しない。

したがって、公理 PP
は\textbf{物理的実現が部分投射像であること}を定める存在公理であり、\textbf{どの像が実現するか}を決める選択力学ではない。

\subsubsection{4.3
物理的対称性}\label{ux7269ux7406ux7684ux5bfeux79f0ux6027}

\textbf{定義 4.1（投射像の物理的対称性）:}
選択像の読出し関係、遷移差集合および合成則を保存する自己同型群を、

\[
G_{\mathrm{phys}}
:=
\operatorname{Aut}
\left(
\operatorname{Im}\mathcal P_*,
\Delta_{\mathcal P_*}
\right)
\]

と定義する。

離散遷移格子で記述する場合は、

\[
G_{\mathrm{lattice}}
:=
\operatorname{Aut}
\left(
\Lambda_{\mathcal P_*},
(\,\cdot\,,\,\cdot\,)
\right)
\]

と書く。ここで格子自己同型は、格子とその内積を同時に保存するものに限る。

物理的対称性は、候補空間全体へ事前に与えられる名称ではなく、選択像に残った関係の保存群として投射後に同定される。

\subsubsection{4.4 非一意分枝}\label{ux975eux4e00ux610fux5206ux679d}

異なる許容部分投射

\[
\mathcal P_1,\mathcal P_2
\in
\mathfrak P(\mathscr S)
\]

が、

\[
\operatorname{Im}\mathcal P_1
\not\cong
\operatorname{Im}\mathcal P_2
\]

を与えることを、本公理は排除しない。したがって、異なる安定条件、読出し条件または境界条件のもとで、異なる対称性分枝が実現し得る。

\begin{center}\rule{0.5\linewidth}{0.5pt}\end{center}

\subsection{5.
八成分二クァルテット分枝}\label{ux516bux6210ux5206ux4e8cux30afux30a1ux30ebux30c6ux30c3ux30c8ux5206ux679d}

以下では、公理 PP が許す分枝の一つに追加条件を置く。これらは公理 PP
の帰結ではなく、\(E_8\) 分枝を定める明示的な分枝仮定である。

\subsubsection{5.1 分枝仮定
Q1：正定値八成分読出し}\label{ux5206ux679dux4eeeux5b9a-q1ux6b63ux5b9aux5024ux516bux6210ux5206ux8aadux51faux3057}

選択像の遷移差が、正定値内積を持つ八次元実読出し空間

\[
V_{\mathrm{read}}
\cong
\mathbb R^8
\]

に入り、

\[
\operatorname{rank}
\Lambda_{\mathcal P_*}
{}=8
\]

を満たすと仮定する。

公理1の複素双線形形式

\[
Q(X)=\sum_jx_j^2
\]

と、ここで導入する正定値読出し内積

\[
(\alpha,\beta)
\]

は別の構造である。\(Q(X)=0\) から正定値性を直接導いてはいない。

\subsubsection{5.2 分枝仮定
Q2：二クァルテット分解}\label{ux5206ux679dux4eeeux5b9a-q2ux4e8cux30afux30a1ux30ebux30c6ux30c3ux30c8ux5206ux89e3}

読出し空間が、二つの四成分部分空間へ直交分解されると仮定する。

\[
V_{\mathrm{read}}
{}=
V_A\oplus V_B,
\qquad
\dim V_A
{}=
\dim V_B
{}=4.
\]

この段階では、各成分に空間、時間、質量、電荷などの物理名を与えない。

\subsubsection{\texorpdfstring{5.3 分枝仮定 Q3：局所 \(D_4\)
遷移格子}{5.3 分枝仮定 Q3：局所 D\_4 遷移格子}}\label{ux5206ux679dux4eeeux5b9a-q3ux5c40ux6240-d_4-ux9077ux79fbux683cux5b50}

各四成分系の局所遷移格子が、座標条件

\[
D_4
:=
\left\{
z\in\mathbb Z^4
\ \middle|\
\sum_{i=1}^{4}z_i
\equiv0
\pmod2
\right\}
\]

で定まる格子と同型であると仮定する。この格子を、数学的分類名として
\(D_4\)
と書く。分枝仮定の内容は上の座標条件であり、対称群名を理由として格子を選んでいるのではない。

したがって、接着前の局所格子は、

\[
L_0
:=
D_4^{(A)}
\oplus
D_4^{(B)}.
\]

この \(D_4\) 型局所構造は、基礎公理または公理 PP
だけからは未導出である。

\subsubsection{\texorpdfstring{5.4 \(D_4\)
の判別類}{5.4 D\_4 の判別類}}\label{d_4-ux306eux5224ux5225ux985e}

\(D_4\) の双対格子を \(D_4^*\) とする。判別群は、

\[
\mathcal A
:=
D_4^*/D_4
\cong
(\mathbb Z/2\mathbb Z)^2
\]

であり、四つの類を、

\[
\mathcal A
{}=
\{0,v,s,c\}
\]

と書く。

代表元は、標準正規直交基底 \(e_1,\ldots,e_4\) を用いて、

\[
v=e_1+D_4,
\]

\[
s=
\frac12
(e_1+e_2+e_3+e_4)
{}+D_4,
\]

\[
c=
\frac12
(-e_1+e_2+e_3+e_4)
{}+D_4
\]

と取れる。

判別二次形式

\[
q:
\mathcal A
\longrightarrow
\mathbb Q/2\mathbb Z,
\qquad
q(a)
{}=
(a,a)
\bmod
2\mathbb Z
\]

は、

\[
q(v)
{}=
q(s)
{}=
q(c)
{}=1
\pmod{2\mathbb Z}
\]

を満たす。

\begin{center}\rule{0.5\linewidth}{0.5pt}\end{center}

\subsection{6.
共通中心中性射影}\label{ux5171ux901aux4e2dux5fc3ux4e2dux6027ux5c04ux5f71}

\subsubsection{6.1 分枝仮定
Q4：共通中心作用}\label{ux5206ux679dux4eeeux5b9a-q4ux5171ux901aux4e2dux5fc3ux4f5cux7528}

二つの四成分系の判別群を独立なラベル群として残さず、単一の共通中心ラベル群

\[
\mathcal A
\cong
(\mathbb Z/2\mathbb Z)^2
\]

へ同定し、同じ元 \(z\in\mathcal A\)
の作用を受けると仮定する。ここで用いるのは、\(\mathcal A\)
の有限アーベル群構造と判別ペアリングである。

判別形式から得られる非退化ペアリングにより、\(\mathcal A\) とその指標群
\(\widehat{\mathcal A}\) を同定する。類 \(a\in\mathcal A\)
に対応する中心指標を \(\chi_a\) とする。\(\mathcal A\)
の指標はすべて実数値 \(\pm1\) を取る。二つの四成分系の判別類を、

\[
(a,b)
\in
\mathcal A\oplus\mathcal A
\]

と書く。

\subsubsection{6.2 分枝仮定
Q5：中性成分の選択}\label{ux5206ux679dux4eeeux5b9a-q5ux4e2dux6027ux6210ux5206ux306eux9078ux629e}

共通中心作用に対して不変な成分だけが物理的読出し像に残ると仮定する。中性射影を、

\[
\Pi_{\mathrm{com}}(a,b)
:=
\frac{1}{|\mathcal A|}
\sum_{z\in\mathcal A}
\chi_a(z)\chi_b(z)
\]

とする。

有限アーベル群の指標直交関係から、

\[
\Pi_{\mathrm{com}}(a,b)
{}=
\delta_{a+b,\,0}.
\]

\subsubsection{6.3
対角類選択定理}\label{ux5bfeux89d2ux985eux9078ux629eux5b9aux7406}

\textbf{定理 6.1（共通中心による対角類選択）:} 分枝仮定 Q4・Q5
のもとで、共通中心中性射影を通過する判別類は、

\[
\boxed{
(0,0),
\qquad
(v,v),
\qquad
(s,s),
\qquad
(c,c)
}
\]

に限られる。

\textbf{証明:} 中性条件は、

\[
a+b=0
\]

である。\(\mathcal A\cong(\mathbb Z/2\mathbb Z)^2\)
の各元は位数1または2なので、

\[
{}-a=a.
\]

したがって、

\[
b=-a=a.
\]

\(\mathcal A=\{0,v,s,c\}\) を代入すれば、上の四類だけが残る。
\(\square\)

\subsubsection{6.4 対角接着群}\label{ux5bfeux89d2ux63a5ux7740ux7fa4}

許容接着群を、

\[
H_{\mathrm{diag}}
:=
\left\{
(0,0),
(v,v),
(s,s),
(c,c)
\right\}
\]

とする。

片側の \(v,s,c\) のラベルを triality
で一括して取り替えたグラフ接着は、判別類の同定を取り直せば同じ対角形で表せる。したがって「同じ類どうし」という記述は、二つの判別群を同一の物理的中心ラベルで同定した後の標準形である。

\begin{center}\rule{0.5\linewidth}{0.5pt}\end{center}

\subsection{\texorpdfstring{7. \(E_8\)
格子の条件付き同定}{7. E\_8 格子の条件付き同定}}\label{e_8-ux683cux5b50ux306eux6761ux4ef6ux4ed8ux304dux540cux5b9a}

\subsubsection{7.1 接着格子}\label{ux63a5ux7740ux683cux5b50}

対角接着後の格子を、

\[
\Lambda
:=
\bigcup_{h\in H_{\mathrm{diag}}}
\left(
L_0+h
\right)
\]

と定義する。

格子の接着と判別形式の一般論は
Nikulin［2］に従う。ただし、以下で必要な指数、行列式および偶性は直接計算する。

\subsubsection{7.2
指数と行列式}\label{ux6307ux6570ux3068ux884cux5217ux5f0f}

\(D_4\) の行列式は4なので、

\[
\det L_0
{}=
\det(D_4\oplus D_4)
{}=
4^2
{}=16.
\]

\(H_{\mathrm{diag}}\) の位数は4であるから、

\[
[\Lambda:L_0]
{}=4.
\]

有限指数部分格子の行列式公式から、

\[
\det\Lambda
{}=
\frac{\det L_0}
{[\Lambda:L_0]^2}
{}=
\frac{16}{4^2}
{}=1.
\]

したがって \(\Lambda\) はユニモジュラーである。

\subsubsection{7.3 偶性}\label{ux5076ux6027}

直和判別形式を \(q\oplus q\) とする。各非零対角類に対して、

\[
(q\oplus q)(a,a)
{}=
q(a)+q(a)
{}=
1+1
{}=0
\pmod{2\mathbb Z},
\]

\[
a\in\{v,s,c\}.
\]

よって \(H_{\mathrm{diag}}\) は判別形式の等方部分群であり、接着後の
\(\Lambda\) は偶格子である。

\subsubsection{7.4 主定理}\label{ux4e3bux5b9aux7406}

\textbf{定理 7.1（八成分二クァルテット分枝の \(E_8\) 同定）:} 分枝仮定
Q1 から Q5 のもとで、共通中心中性射影によって得られる遷移差格子
\(\Lambda\) は \(E_8\) 格子と同型である。

\[
\boxed{
\Lambda
\cong
\Lambda_{E_8}.
}
\]

\textbf{証明:} Q1・Q2により \(\Lambda\)
は正定値八次元である。Q3により接着前の格子は \(D_4\oplus D_4\)
である。定理6.1により接着群は \(H_{\mathrm{diag}}\) となる。§7.2から
\(\det\Lambda=1\)、§7.3から \(\Lambda\) は偶である。したがって
\(\Lambda\)
は正定値・八次元・偶・ユニモジュラー格子である。八次元におけるこの格子は同型を除いて
\(E_8\) 格子に一意である［3］。ゆえに \(\Lambda\cong\Lambda_{E_8}\)。
\(\square\)

\subsubsection{7.5 論理的身分}\label{ux8ad6ux7406ux7684ux8eabux5206}

定理7.1は、

\[
Q(X)=0
\Longrightarrow
E_8
\]

という無条件命題ではない。正確な主張は、

\[
\boxed{
\begin{array}{c}
\text{公理 PP}\\
{}+\text{Q1：正定値八成分}\\
{}+\text{Q2：二クァルテット}\\
{}+\text{Q3：局所 }D_4\\
{}+\text{Q4・Q5：共通中心中性射影}
\end{array}
\Longrightarrow
\Lambda_{E_8}.
}
\]

したがって \(E_8\)
は、公理に埋め込まれた名称ではなく、条件を満たす投射像の分類結果である。

\begin{center}\rule{0.5\linewidth}{0.5pt}\end{center}

\subsection{8. 240根と248次元}\label{ux6839ux3068248ux6b21ux5143}

\subsubsection{8.1
最短遷移ベクトル}\label{ux6700ux77edux9077ux79fbux30d9ux30afux30c8ux30eb}

\(E_8\) 格子のノルム2ベクトルを、

\[
\Phi
:=
\left\{
\alpha\in\Lambda
\mid
(\alpha,\alpha)=2
\right\}
\]

とする。

\subsubsection{\texorpdfstring{8.2 \(D_4\oplus D_4\)
内部}{8.2 D\_4\textbackslash oplus D\_4 内部}}\label{d_4oplus-d_4-ux5185ux90e8}

\(D_4\) の根は、

\[
\pm e_i\pm e_j,
\qquad
1\le i<j\le4
\]

であり、その本数は、

\[
4
\binom42
{}=24.
\]

したがって、

\[
D_4\oplus D_4
\]

内部のノルム2ベクトルは、

\[
24+24=48
\]

本である。

\subsubsection{8.3 非零接着類}\label{ux975eux96f6ux63a5ux7740ux985e}

\(v,s,c\)
の各判別類は、ノルム1の最短代表を8本ずつ持つ。したがって、各対角類

\[
(v,v),
\qquad
(s,s),
\qquad
(c,c)
\]

は、

\[
8\times8=64
\]

本のノルム2ベクトルを与える。

非零対角類は三つなので、

\[
\boxed{
|\Phi|
{}=
48
{}+3\times64
{}=240.
}
\]

\subsubsection{8.4 鏡映閉包}\label{ux93e1ux6620ux9589ux5305}

各 \(\alpha\in\Phi\) に対し、

\[
s_\alpha(x)
:=
x
{}-\frac{2(x,\alpha)}
{(\alpha,\alpha)}
\alpha
{}=
x-(x,\alpha)\alpha
\]

とする。

\(\Lambda\) は整数格子であり \((\alpha,\alpha)=2\) なので、

\[
s_\alpha(\Lambda)
{}=
\Lambda.
\]

これらの鏡映が生成する群は、\(E_8\) 根系の Weyl 群 \(W(E_8)\) である。

\subsubsection{8.5 Lie 代数}\label{lie-ux4ee3ux6570}

八次元 Cartan 空間 \(\mathfrak h\) と各根に対応する一次元根空間
\(\mathfrak g_\alpha\) を用いて、

\[
\mathfrak e_8(\mathbb C)
{}=
\mathfrak h
\oplus
\bigoplus_{\alpha\in\Phi}
\mathfrak g_\alpha
\]

を得る。

したがって、

\[
\boxed{
\dim\mathfrak e_8
{}=
8+240
{}=248.
}
\]

離散的な遷移格子の鏡映対称性は \(W(E_8)\) であり、根空間を含む複素単純
Lie 代数が \(\mathfrak e_8(\mathbb C)\)
である。正定値内積に対応するコンパクト実形を選ぶ場合、その単連結 Lie
群がコンパクト型 \(E_8\) となる。

\begin{center}\rule{0.5\linewidth}{0.5pt}\end{center}

\subsection{9.
有限回帰との接続}\label{ux6709ux9650ux56deux5e30ux3068ux306eux63a5ux7d9a}

\subsubsection{9.1 Coxeter 元}\label{coxeter-ux5143}

\(E_8\) 根系の Coxeter 数は30であり、その Coxeter 元 \(C\) は、

\[
\boxed{
C^{30}=I
}
\]

を満たす［4］。

Cartan 空間上の固有値は、

\[
\operatorname{Spec}(C)
{}=
\left\{
\exp
\left(
\frac{2\pi i m}{30}
\right)
\ \middle|\
m\in
\{1,7,11,13,17,19,23,29\}
\right\}.
\]

240根は、Coxeter 作用のもとで長さ30の八軌道に分かれる［4］。

\[
\Phi
{}=
\bigsqcup_{a=1}^{8}
\mathcal O_a,
\qquad
|\mathcal O_a|
{}=30.
\]

したがって、

\[
\boxed{
248
{}=
8+240
{}=
8+8\times30
{}=
8(30+1).
}
\]

\subsubsection{9.2
有限回帰公理との論理関係}\label{ux6709ux9650ux56deux5e30ux516cux7406ux3068ux306eux8ad6ux7406ux95a2ux4fc2}

\(E_8\) 投射像には、位数30の有限回帰作用素 \(C\) が存在する。しかし、

\[
U^n=I
\]

を満たす一般の物理的発展作用素 \(U\) が、

\[
U=C
\]

であることは、定理7.1からは導かれない。

したがって、

\[
U_{\mathrm{phys}}
\longleftrightarrow
C
\]

の同定には、物理的発展が根系全体を Coxeter
軌道として巡回させることを示す追加の対応写像が必要である。

本稿が確立するのは、

\[
\boxed{
\Lambda_{\mathcal P_*}
\cong
\Lambda_{E_8}
\Longrightarrow
\exists C\in\operatorname{Aut}(\Lambda_{\mathcal P_*})
\text{ such that }
C^{30}=I
}
\]

までである。

\begin{center}\rule{0.5\linewidth}{0.5pt}\end{center}

\subsection{10.
無名性との整合}\label{ux7121ux540dux6027ux3068ux306eux6574ux5408}

\subsubsection{10.1
対称群名を投射条件に用いない}\label{ux5bfeux79f0ux7fa4ux540dux3092ux6295ux5c04ux6761ux4ef6ux306bux7528ux3044ux306aux3044}

公理 PP と許容性条件は、\(E_8\) という名称を含まない。分枝仮定も、

\begin{enumerate}
\def\labelenumi{\arabic{enumi}.}
\tightlist
\item
  八成分性
\item
  二つの四成分系
\item
  局所遷移差の偶整数条件
\item
  共通中心中性射影
\end{enumerate}

として記述される。

その結果得られた格子を、既知の格子分類によって \(E_8\)
と同定する。したがって、

\[
\boxed{
\text{投射条件}
\longrightarrow
\text{格子構造}
\longrightarrow
\text{対称群名}
}
\]

の順序を守る。

\subsubsection{10.2
物理名も投射後に与える}\label{ux7269ux7406ux540dux3082ux6295ux5c04ux5f8cux306bux4e0eux3048ux308b}

八つの成分を、後に

\[
(x,y,z,R)
\oplus
(t,Q_1,Q_2,Q_3)
\]

などと解釈することは、本稿の定理には含まれない。この対応は、各成分の読出し演算、保存量、交換則および実験対応を定めた後にのみ与えられる物理的解釈である。

\subsubsection{10.3
E8以外の分枝}\label{e8ux4ee5ux5916ux306eux5206ux679d}

分枝仮定 Q1 から Q5 のいずれかが成立しない場合、本稿は \(E_8\)
を要求しない。

\begin{itemize}
\tightlist
\item
  読出しランクが8でなければ、定理7.1の適用外である。
\item
  局所格子が \(D_4\) 型でなければ、別の接着問題になる。
\item
  対角接着がなければ、\(D_4\oplus D_4\)
  は行列式16の分解可能格子のままである。
\item
  正定値性がなければ、正定値 \(E_8\) 格子の一意性定理は適用できない。
\end{itemize}

したがって、本枠組みは別の部分投射像と別の対称性分枝を最初から許している。

\begin{center}\rule{0.5\linewidth}{0.5pt}\end{center}

\subsection{11. 選択問題}\label{ux9078ux629eux554fux984c}

\subsubsection{11.1
未導出の中心}\label{ux672aux5c0eux51faux306eux4e2dux5fc3}

本稿の未導出部分は、格子接着後の計算ではない。未導出なのは、実際の物理系がなぜ分枝仮定
Q1 から Q5 を満たす投射を選ぶかである。

特に必要なのは、

\[
\boxed{
\mathcal P_*
\text{ を選ぶ無名な選択則}
}
\]

である。

\subsubsection{11.2
選択汎関数による表示}\label{ux9078ux629eux6c4eux95a2ux6570ux306bux3088ux308bux8868ux793a}

選択則が汎関数として書ける場合、

\[
\mathcal F:
\mathfrak P(\mathscr S)
\longrightarrow
\mathbb R
\]

を導入し、

\[
\mathcal P_*
\in
\operatorname*{arg\,ext}_{\mathcal P\in\mathfrak P(\mathscr S)}
\mathcal F(\mathcal P)
\]

と書ける。

\(\operatorname*{arg\,ext}\)
は最大、最小または停留条件のいずれかを表す。本稿は \(\mathcal F\)
の具体形を仮定しない。

\subsubsection{11.3
選択則へ要求される条件}\label{ux9078ux629eux5247ux3078ux8981ux6c42ux3055ux308cux308bux6761ux4ef6}

選択則は少なくとも次を満たさなければならない。

\begin{enumerate}
\def\labelenumi{\arabic{enumi}.}
\tightlist
\item
  個別成分名を参照しない。
\item
  絶対スケールを参照しない。
\item
  閉鎖保存を破らない。
\item
  対称群名を目的関数へ直接代入しない。
\item
  八成分性、局所 \(D_4\)
  構造、共通中心作用が選ばれる理由を、同一の基準で比較可能にする。
\end{enumerate}

\subsubsection{11.4
共通中心条件の位置}\label{ux5171ux901aux4e2dux5fc3ux6761ux4ef6ux306eux4f4dux7f6e}

共通中心中性射影は、対角接着を導く十分条件である。しかし、

\[
\text{なぜ同じ中心が二クァルテットへ作用するのか}
\]

および、

\[
\text{なぜ中性成分だけが読出されるのか}
\]

は、選択力学または観測構成から導出されなければならない。

したがって共通中心条件は、現時点では \(E_8\)
分枝の\textbf{選択公理候補}であり、基礎公理からの導出済み帰結ではない。

\begin{center}\rule{0.5\linewidth}{0.5pt}\end{center}

\subsection{12. 判別条件}\label{ux5224ux5225ux6761ux4ef6}

本枠組みは、次の判別計算を要求する。

\subsubsection{12.1
読出しランク}\label{ux8aadux51faux3057ux30e9ux30f3ux30af}

物理的に独立な遷移差のランクが、

\[
\operatorname{rank}\Lambda_{\mathcal P_*}
{}=8
\]

であるかを検査する。

\subsubsection{12.2 局所格子}\label{ux5c40ux6240ux683cux5b50}

各四成分系の最小遷移ベクトルが、\(D_4\) の根条件

\[
\pm e_i\pm e_j
\]

で閉じるかを検査する。

\subsubsection{12.3 判別類相関}\label{ux5224ux5225ux985eux76f8ux95a2}

二つの四成分系において、判別類の同時出現が、

\[
(a,b)
\text{ が出現}
\quad\Longleftrightarrow\quad
a=b
\]

を満たすかを検査する。

非対角類

\[
(v,s),
\quad
(v,c),
\quad
(s,c)
\]

が安定に残るなら、共通中心中性射影仮説は棄却される。

\subsubsection{12.4 Coxeter 周期}\label{coxeter-ux5468ux671f}

選択された遷移差全体が、八つの長さ30軌道へ分かれるかを検査する。

\[
240
\stackrel{?}{=}
8\times30.
\]

この検査が成立しても、物理的時間発展 \(U\) と Coxeter 元 \(C\)
の同一性には、作用素レベルの対応写像が別途必要である。

\begin{center}\rule{0.5\linewidth}{0.5pt}\end{center}

\subsection{13.
基本公理系への編入形}\label{ux57faux672cux516cux7406ux7cfbux3078ux306eux7de8ux5165ux5f62}

本稿の成果を基本公理系へ編入する場合、追加すべき一般原理は \(E_8\)
そのものではなく、公理 PP である。

最小編入形は次である。

\subsubsection{13.1 編入公理案}\label{ux7de8ux5165ux516cux7406ux6848}

\begin{quote}
\textbf{無名部分投射存在公理:}
基礎公理の許容解空間に対し、物理的に実現される状態空間は、無名性、スケール無名性、閉鎖整合性および読出し整合性を満たす少なくとも一つの非自明な許容部分投射の像である。
\end{quote}

\[
\mathscr S_{\mathrm{phys}}
{}=
\operatorname{Im}\mathcal P_*,
\qquad
\mathcal P_*
\in
\mathfrak P(\mathscr S).
\]

\subsubsection{13.2
編入しないもの}\label{ux7de8ux5165ux3057ux306aux3044ux3082ux306e}

次は基礎公理本文へ直接編入しない。

\begin{enumerate}
\def\labelenumi{\arabic{enumi}.}
\tightlist
\item
  \(E_8\) という対称群名。
\item
  240、248、30という導出値。
\item
  八成分二クァルテット分枝だけを唯一の実現とする宣言。
\end{enumerate}

これらは、公理 PP
に追加分枝条件を置いた本稿の定理および帰結として引用する。

\subsubsection{13.3
単一正本との整合}\label{ux5358ux4e00ux6b63ux672cux3068ux306eux6574ux5408}

公開済みの基本公理系 v8［1］を改変せず、本稿の DOI
公開後、v8を基礎とする次版に公理 PP を追加する。これにより、

\[
\text{独立論文で定義と定理を固定}
\longrightarrow
\text{基本公理系次版へ一般公理だけを編入}
\]

という順序を保つ。

\begin{center}\rule{0.5\linewidth}{0.5pt}\end{center}

\subsection{14. 結論}\label{ux7d50ux8ad6}

本稿は、物理的対称性を基礎公理の全解空間へ事前に与えるのではなく、許容部分投射像の自己同型として投射後に同定する公理的枠組みを提示した。

中心公理は、

\[
\boxed{
\mathscr S_{\mathrm{phys}}
{}=
\operatorname{Im}\mathcal P_*,
\qquad
\mathcal P_*
\in
\mathfrak P(\mathscr S)
}
\]

である。

この公理は、どの対称群を選ぶかを指定しない。成分名、絶対スケール、物理名、対称群名に依存せず、閉鎖と読出しに整合する部分投射像が物理的実現を与える、とだけ定める。

その特定分枝として、正定値八成分読出し、二クァルテット分解、局所 \(D_4\)
遷移格子、共通中心中性射影を仮定した。このとき、

\[
\boxed{
(0,0),
\qquad
(v,v),
\qquad
(s,s),
\qquad
(c,c)
}
\]

だけが残り、接着格子は、

\[
\boxed{
\Lambda
\cong
\Lambda_{E_8}
}
\]

となる。

そこから、

\[
\boxed{
|\Phi|
{}=240,
\qquad
\dim\mathfrak e_8
{}=248,
\qquad
C^{30}
{}=I
}
\]

が得られる。

したがって、本稿における \(E_8\) の論理的位置は、

\[
\boxed{
\text{$E_8$ を選んだのではなく、}
\quad
\text{無名な部分投射が選んだ像を分類した結果が $E_8$ である}
}
\]

である。

残る中心問題は、許容部分投射の集合から実際の \(\mathcal P_*\)
を選ぶ無名な選択汎関数または安定性条件の導出である。

\begin{center}\rule{0.5\linewidth}{0.5pt}\end{center}

\subsection{参考文献}\label{ux53c2ux8003ux6587ux732e}

［1］木原 範昭, 「無名等振幅複合波モデル基本公理系 v8」, Zenodo, Version
DOI:
\href{https://doi.org/10.5281/zenodo.21495422}{10.5281/zenodo.21495422},
Concept DOI:
\href{https://doi.org/10.5281/zenodo.21315735}{10.5281/zenodo.21315735},
2026.

［2］V. V. Nikulin, ``Integral symmetric bilinear forms and some of
their applications,'' \emph{Mathematics of the USSR-Izvestiya}, vol.~14,
no. 1, pp.~103--167, 1980. DOI:
\href{https://doi.org/10.1070/IM1980v014n01ABEH001060}{10.1070/IM1980v014n01ABEH001060}.

［3］L. J. Mordell, ``The definite quadratic forms in eight variables
with determinant unity,'' \emph{Journal de Mathématiques Pures et
Appliquées}, 9e série, vol.~17, pp.~41--46, 1938.
\href{https://www.numdam.org/item/JMPA_1938_9_17_1-4_41_0/}{NUMDAM}.

［4］David A. Richter, ``Triacontagonal coordinates for the \(E_8\) root
system,'' arXiv:0704.3091, 2007.
\href{https://arxiv.org/abs/0704.3091}{arXiv}.

\begin{center}\rule{0.5\linewidth}{0.5pt}\end{center}

\subsection{付録A：主張の身分監査}\label{ux4ed8ux9332aux4e3bux5f35ux306eux8eabux5206ux76e3ux67fb}

{\def\LTcaptype{none} % do not increment counter
\begin{longtable}[]{@{}
  >{\raggedright\arraybackslash}p{(\linewidth - 4\tabcolsep) * \real{0.3333}}
  >{\raggedright\arraybackslash}p{(\linewidth - 4\tabcolsep) * \real{0.3333}}
  >{\raggedright\arraybackslash}p{(\linewidth - 4\tabcolsep) * \real{0.3333}}@{}}
\toprule\noalign{}
\begin{minipage}[b]{\linewidth}\raggedright
対象
\end{minipage} & \begin{minipage}[b]{\linewidth}\raggedright
身分
\end{minipage} & \begin{minipage}[b]{\linewidth}\raggedright
本稿での扱い
\end{minipage} \\
\midrule\noalign{}
\endhead
\bottomrule\noalign{}
\endlastfoot
無名性・スケール無名性・零閉鎖 & 既存公理 & 基本公理系
v8［1］から継承 \\
物理的実現が許容部分投射像である & 新公理 PP & 本稿の中心公理 \\
許容部分投射が複数存在し得る & 公理 PP が許す構造 &
一意性を仮定しない \\
実際の \(\mathcal P_*\) の選択則 & 未導出 &
選択汎関数または力学が必要 \\
正定値八成分読出し & 分枝仮定 Q1 & 基礎公理から未導出 \\
二クァルテット分解 & 分枝仮定 Q2 & 基礎公理から未導出 \\
局所 \(D_4\) 遷移格子 & 分枝仮定 Q3 & 基礎公理から未導出 \\
共通中心作用 & 分枝仮定 Q4 & 選択公理候補 \\
中性成分だけを残す & 分枝仮定 Q5 & 選択公理候補 \\
対角類だけが残る & 定理6.1 & Q4・Q5から導出 \\
偶・ユニモジュラー八次元格子 & 導出済み帰結 &
接着指数と判別形式から導出 \\
\(\Lambda\cong\Lambda_{E_8}\) & 条件付き定理7.1 & Q1--Q5のもとで成立 \\
240根 & 導出済み帰結 & 接着類ごとの直接計数 \\
248次元 & 導出済み帰結 & rank 8と240根から導出 \\
\(C^{30}=I\) & \(E_8\) 同定後の既知数学的帰結 & Richter［4］ \\
物理的 \(U\) と Coxeter 元 \(C\) の同一視 & 接続課題 &
作用素対応写像が必要 \\
八成分と \(xyzR,tQ_1Q_2Q_3\) の対応 & 物理的解釈・未導出 &
本稿の定理に含めない \\
\end{longtable}
}

\begin{center}\rule{0.5\linewidth}{0.5pt}\end{center}

\subsection{付録B：分岐監査}\label{ux4ed8ux9332bux5206ux5c90ux76e3ux67fb}

{\def\LTcaptype{none} % do not increment counter
\begin{longtable}[]{@{}
  >{\raggedright\arraybackslash}p{(\linewidth - 2\tabcolsep) * \real{0.5000}}
  >{\raggedright\arraybackslash}p{(\linewidth - 2\tabcolsep) * \real{0.5000}}@{}}
\toprule\noalign{}
\begin{minipage}[b]{\linewidth}\raggedright
条件を外した場合
\end{minipage} & \begin{minipage}[b]{\linewidth}\raggedright
結果
\end{minipage} \\
\midrule\noalign{}
\endhead
\bottomrule\noalign{}
\endlastfoot
公理 PPを外す & 候補空間と物理的実現空間を区別する原理がなくなる \\
Q1を外す & 正定値八次元格子の同定定理を適用できない \\
Q2を外す & \(D_4\oplus D_4\) 接着として記述できない \\
Q3を外す & 判別群 \(\{0,v,s,c\}\) と対角接着の起点が失われる \\
Q4を外す & 二クァルテットの判別類を独立に選べる \\
Q5を外す & 非対角類を除く射影根拠がなくなる \\
対角接着を外す & \(D_4\oplus D_4\) は行列式16で自己双対でない \\
正定値性を不定符号へ変える & 正定値 \(E_8\)
格子の一意性は適用外となる \\
物理的 \(U=C\) を仮定しない & \(E_8\)
内の位数30作用は存在するが、物理時間発展との同一性は未確定 \\
\end{longtable}
}

\begin{center}\rule{0.5\linewidth}{0.5pt}\end{center}

\subsection{付録C：最小導出鎖}\label{ux4ed8ux9332cux6700ux5c0fux5c0eux51faux9396}

\[
\boxed{
\begin{aligned}
&\text{無名性・スケール無名性・零閉鎖}\\
&\Longrightarrow
\mathscr S\\
&\xrightarrow{\text{公理 PP}}
\operatorname{Im}\mathcal P_*\\
&\xrightarrow{\text{Q1--Q3}}
D_4\oplus D_4\\
&\xrightarrow{\text{Q4--Q5}}
H_{\mathrm{diag}}
{}=
\{(0,0),(v,v),(s,s),(c,c)\}\\
&\Longrightarrow
\text{正定値八次元偶ユニモジュラー格子}\\
&\Longrightarrow
\Lambda_{E_8}\\
&\Longrightarrow
240\text{ roots}\\
&\Longrightarrow
\dim\mathfrak e_8
{}=248\\
&\Longrightarrow
\exists C,\quad C^{30}=I.
\end{aligned}
}
\]

\end{document}
